Two-stage engineering architecture of HD-MCRA.
Stage 1 trains and freezes a velocity tracker, then uses paired nominal and executed command rollouts to calibrate cost and trajectory margins; $Q$ denotes the chosen upper empirical quantile for each error.
The safety scene includes an equal-scale local enlargement of the paired paths; its orange segment connects nominal and executed positions at the same time, and that individual deviation is distinct from the calibrated margin $\delta$.
Stage 2 uses the same fixed margins: $\tilde c=(1+\alpha)c$ supplies both cost-learning targets and budget updates, while $\mathcal F_\delta=\{x:d(x,\mathcal F)\leq\delta\}$ expands the unsafe region before forming reachability targets.
The parallel cost and reachability critics support the joint actor update and budget-feasibility evaluation; their dashed branches merge into one learning link to the actor.
The actor sends high-level velocity commands through the same frozen tracker, which generates low-level joint commands for execution.
Two compact scene insets, with hierarchy labels, are centered immediately above their corresponding execution arrows: body translation and yaw arrows illustrate the actor's velocity command, while curved arrows at the front-leg joints illustrate the tracker's joint-level commands.
These local arrows indicate command categories, not measured velocities, joint targets, or torques.
Solid architecture arrows show the module input and execution chain.
Three Go2 scene views represent navigation observations, the goal location, and remaining cost budget $z_t$; the route in the budget view illustrates task progress, not battery charge.
The terminal interaction inset combines three static placements of the same robot at different times along the existing illustrative executed path, with increasing opacity, three obstacles, and the goal.
Ghost poses indicate time progression rather than multiple robots.
Observation fan-out and execution feedback connections are omitted for clarity; the controller still operates in closed loop.
Network nodes and layers are schematic.
All scene views are static renders, not camera measurements or policy rollouts; displayed paths and underlying example costs are assumed illustrations, not measured calibration statistics.
